\section{Method}
\label{sec:method}

\subsection{Background}

Given queries $Q$, keys $K$, and values $V$ in $\mathbb{R}^{n \times d}$, scaled dot-product attention computes
\begin{equation}
  \operatorname{Attn}(Q, K, V) = \softmax\!\left( \frac{Q K^\top}{\sqrt{d}} \right) V ,
  \label{eq:attention}
\end{equation}
which needs $O(n^2)$ time and memory.

\subsection{Block summaries}

We partition the sequence into $B = \lceil n / b \rceil$ contiguous blocks of $b$ tokens and summarize block $j$ by the mean of its keys,
\begin{equation}
  \bar{k}_j = \frac{1}{b} \sum_{i \in \mathcal{B}_j} k_i .
\end{equation}
Each query $q_i$ scores every block by $s_{ij} = q_i^\top \bar{k}_j / \sqrt{d}$ and keeps the $r$ highest-scoring blocks. The token then attends to the tokens of its local window of width $w$ and to the summaries of the selected blocks.

\subsection{Complexity}

Scoring costs $O(nB)$, and attention costs $O(n(w + r))$. With $b = \sqrt{n}$ the total is $O(n\sqrt{n})$; with a fixed number of blocks it is $O(n)$. In all experiments we use $w = 256$, $b = 64$, and $r = 8$.

\subsection{Training}

Hard top-$r$ selection has zero gradient almost everywhere. During training we relax it with the perturbed-maximum estimator, adding Gumbel noise to the scores and annealing the temperature from $1.0$ to $0.1$ over the first 10{,}000 steps.
